\scanpage{100}
indexée par $\underline\omega\times\underline t^*$ soit une famille d'éléments de $V_I(k)$ s'exprime par la formule
\[
\sum_{i\in I}\psi_i^\omega(\delta)=0 \tag*{(75)}
\]
les \deleted{\unsure{f}} conditions (71) et (71) s'expriment aussitôt en termes des $\psi_i^\omega(\delta)$, comme
\[
\psi_i^{\omega'}(\delta)=i_\omega\psi_i^\omega(\delta) \tag*{(76)}
\]
\[
\psi_i^\omega(\lambda\delta)=\lambda^{-1}\psi_i^\omega(\delta)\,. \tag*{(77)}
\]
Il reste à exprimer le fait que les éléments $\psi^\omega(\delta)\in V_I(k)$ définis par les $(\omega,\delta)\in\underline\omega\times\underline t\deleted{\unsure{_E}}^*$ sont dans $F_t$ -- moyennant (76) (77) il suffit de l'exprimer pour un seul couple $(\omega,\delta)$ --, et la condition (73). \deleted{On exprime les deux simultanément, en écrivant}
\inserted{Si on définit $F_t$ comme noyau d'une forme linéaire sur $V_I$}
\[
\deleted{\sum_{i\in I}\psi_i^\omega(\delta)e_i-}\tag*{\deleted{(77)}}
\]
\[
F_t=\deleted{\unsure{\sum}}\left\{\sum\lambda_ie_i\mathrel{}\middle|\mathrel{}\sum\lambda_i=0,\ 
\sum c_i\lambda_i\deleted{\unsure{\psi_i^\omega(\delta)|e_i}}=0\right\}
\]
\margin{$(e_i)_{i\in I}$ base canonique de $k^I$}
[\,\uline{NB} l'élément $(c_i)_{i\in I}$ de $\deleted{\unsure{\Sigma}}k^I$ est déterminé modulo multiplication par $k^*$ et addition d'un élément diagonal\,], \deleted{cette}la condition $\psi_{\unsure{\cdot|i}}^\omega(\delta)\in F_t$ s'exprime par la condition calculatoire
\[
\sum_{i\in I}c_i\psi_i^\omega(\delta)=0\qquad . \tag*{(78)}
\]
Il reste \inserted{ensuite} à exprimer (73)
% archive sheet number: 99
